\documentclass[11pt,a4paper]{article}

\usepackage[utf8]{inputenc}
\usepackage[T1]{fontenc}
\usepackage[spanish,es-nodecimaldot]{babel}
\usepackage{geometry}
\usepackage{booktabs}
\usepackage{array}
\usepackage{longtable}
\usepackage{listings}
\usepackage{xcolor}
\usepackage{hyperref}
\usepackage{enumitem}
\usepackage{amsmath}

\geometry{margin=2.5cm}

\hypersetup{
colorlinks=true,
linkcolor=blue,
urlcolor=blue,
citecolor=blue
}

\lstset{
basicstyle=\ttfamily\small,
columns=fullflexible,
keepspaces=true,
frame=single,
breaklines=true,
showstringspaces=false
}

\title{Multi-Push Heap Characterization\
Experimental Runbook}
\author{}
\date{\today}

\begin{document}

\maketitle

\section{Runbook}

This runbook defines the procedure used to characterize the reference
implementation of the Multi-Push optimization and to verify the behavior
of the two state variables that are central to the proposed architecture: push location and heap size.

The experiment is intentionally performed at the software-reference level
before moving to RTL. The objective is not to measure hardware performance.
The objective is to establish a deterministic behavioral model that can
later be used as the reference for an RTL implementation.

\subsection{1. Experimental Objective}

The experiment addresses the following question:

\begin{quote}
When multiple PUSH operations arrive before heap maintenance is completed,
can the Multi-Push mechanism accept those entries immediately, defer their
HeapifyUp operations, and subsequently synchronize the heap before a POP
without changing the priority-queue output sequence?
\end{quote}

The characterization experiment focuses specifically on the relationship

push_location>heap_size

as the software representation of pending heap-maintenance work.

In this model, \texttt{push_location} identifies the first free tail
position for an incoming PUSH, while \texttt{heap_size} identifies the
number of entries that have already been incorporated into the heap
ordering.

Therefore, the difference

pending=push_location−heap_size

represents the number of accepted PUSH operations that have not yet
completed HeapifyUp.

\subsection{2. Software Under Test}

The experiment uses the following reference programs:

\begin{itemize}
\item \texttt{reference/multipush_reference.cpp}: functional
Multi-Push reference model with operation and comparison statistics.

\item \texttt{reference/multipush\_characterize.cpp}: state
characterization model used to expose the relationship between
accepted PUSH operations and completed heap maintenance.

\item \texttt{reference/heap\_drain}: conventional baseline priority
queue used to verify functional output equivalence.

\item \texttt{reference/heap\_phase\_stats}: instrumented baseline
used to obtain PUSH and POP comparison and swap counts.


\end{itemize}

The \texttt{multipush_reference} implementation is kept as the current
functional reference model.

The characterization program is an auxiliary experiment used to inspect
the internal state transition implied by the Multi-Push algorithm. It is
not intended to replace the functional reference model.

\subsection{3. Build and Environment Check}

Start from the repository root:

\begin{lstlisting}
/foss/designs/priority-queue-sky130
\end{lstlisting}

Enter the C++ reference-model directory:

\begin{lstlisting}
cd /foss/designs/priority-queue-sky130/cpp
\end{lstlisting}

Confirm that the reference executables exist:

\begin{lstlisting}
ls -lh reference/
\end{lstlisting}

The relevant executables are:

\begin{lstlisting}
reference/heap_drain
reference/heap_phase_stats
reference/multipush_reference
reference/multipush_characterize
\end{lstlisting}

If a source file has changed, rebuild the corresponding executable before
collecting new results.

The source revision should also be recorded before a result is considered
reproducible:

\begin{lstlisting}
git status --short
git log -1 --oneline
\end{lstlisting}

\subsection{4. Functional Baseline Check}

The first check is functional equivalence against the conventional heap
implementation.

For the generated workload with N=128 and seed 12345, run:

\begin{lstlisting}
./reference/multipush_reference
workloads/generated/workload_128_seed12345.txt
> /tmp/multipush_128.txt
2> /tmp/multipush_128.stats

./reference/heap_drain
workloads/generated/workload_128_seed12345.txt
> /tmp/baseline_128.txt
\end{lstlisting}

Compare the generated POP sequence:

\begin{lstlisting}
diff -u /tmp/baseline_128.txt /tmp/multipush_128.txt
\end{lstlisting}

A clean result means that no differences were reported.

For an explicit pass/fail check:

\begin{lstlisting}
cmp -s /tmp/baseline_128.txt /tmp/multipush_128.txt
&& echo "PASS: N=128 outputs identical"
|| echo "FAIL: N=128 outputs differ"
\end{lstlisting}

The observed result for the current implementation was:

\begin{lstlisting}
PASS: N=128 outputs identical
\end{lstlisting}

This establishes functional equivalence for that workload. It does not
by itself establish equivalence for all possible workloads.

\subsection{5. Collect Multi-Push Statistics}

The Multi-Push reference model reports separate statistics for the PUSH
and POP phases.

Run:

\begin{lstlisting}
cat /tmp/multipush_128.stats
\end{lstlisting}

The current N=128 result is:

\begin{lstlisting}
pushes=128
pops=128
push_comparisons=255
push_swaps=135
pop_comparisons=1157
pop_swaps=562
total_comparisons=1412
total_swaps=697
heapify_up_count=128
\end{lstlisting}

These counters characterize the algorithmic work performed by the
software reference model. They should not be interpreted as direct
hardware cycle counts.

\subsection{6. Collect Baseline Phase Statistics}

The conventional implementation is instrumented separately using
\texttt{heap_phase_stats}.

Run:

\begin{lstlisting}
for N in 8 16 32 64 128; do
./reference/heap_phase_stats
"workloads/generated/workload_${N}seed12345.txt"
> "/tmp/baseline_stats${N}.txt"

echo "===== BASELINE STATS N=${N} ====="
cat "/tmp/baseline_stats_${N}.txt"


done
\end{lstlisting}

The observed baseline results are:

\begin{lstlisting}
N=8:
push_comparisons=11
push_swaps=8
pop_comparisons=16
pop_swaps=8
total_comparisons=27
total_swaps=16

N=16:
push_comparisons=21
push_swaps=10
pop_comparisons=56
pop_swaps=28
total_comparisons=77
total_swaps=38

N=32:
push_comparisons=55
push_swaps=30
pop_comparisons=175
pop_swaps=84
total_comparisons=230
total_swaps=114

N=64:
push_comparisons=116
push_swaps=59
pop_comparisons=452
pop_swaps=223
total_comparisons=568
total_swaps=282

N=128:
push_comparisons=255
push_swaps=135
pop_comparisons=1157
pop_swaps=562
total_comparisons=1412
total_swaps=697
\end{lstlisting}

These values provide a useful sanity check because the current Multi-Push
reference produces the same aggregate operation counts for the
corresponding PUSH-only workloads.

\subsection{7. Directed Multi-Push Characterization}

The most important experiment is the directed characterization workload.

Create the workload:

\begin{lstlisting}
cat > /tmp/multipush_directed.txt <<'EOF'
PUSH 10 10
PUSH 25 25
PUSH 5 5
PUSH 40 40
PUSH 15 15
POP
POP
POP
POP
POP
EOF
\end{lstlisting}

Run:

\begin{lstlisting}
./reference/multipush_characterize
/tmp/multipush_directed.txt
> /tmp/multipush_directed.out
\end{lstlisting}

Inspect the trace:

\begin{lstlisting}
cat /tmp/multipush_directed.out
\end{lstlisting}

The important portion of the trace is:

\begin{lstlisting}
0 PUSH 1 0 1
1 PUSH 2 0 2
2 PUSH 3 0 3
3 PUSH 4 0 4
4 PUSH 5 0 5
5 POP_WAIT 5 0 5
6 POP_READY 5 5 0
7 POP_DONE 4 4 0
\end{lstlisting}

The interpretation is:

\begin{itemize}
\item Five PUSH operations are accepted while
\texttt{heap_size} remains zero.

\item \texttt{push\_location} advances from zero to five.

\item Therefore, five entries are pending heap maintenance.

\item The POP reaches a synchronization point.

\item The pending entries are incorporated before the POP becomes
ready.

\item After the root is removed, the active heap contains four entries.


\end{itemize}

This is the key behavioral observation of the experiment.

\subsection{8. Interleaved Workload}

The second directed test checks that the same mechanism works after the
queue has already been partially maintained.

Run:

\begin{lstlisting}
./reference/multipush_characterize
/tmp/multipush_interleaved.txt
> /tmp/multipush_interleaved.out

cat /tmp/multipush_interleaved.out
\end{lstlisting}

The observed trace contains transitions of the form:

\begin{lstlisting}
6 PUSH 3 2 1
7 PUSH 4 2 2
8 POP_WAIT 4 2 2
9 POP_READY 4 4 0
10 POP_DONE 3 3 0
\end{lstlisting}

This is important because the experiment is no longer starting from an
empty heap.

At cycle 6, the queue already contains two maintained entries. Two
additional PUSH operations are then accepted without immediately
increasing \texttt{heap_size}. The resulting state is

push_location=4,heap_size=2,pending=2.

The following POP forces synchronization. The state changes to

push_location=4,heap_size=4,pending=0.

Only after this synchronization does the root removal occur.

\subsection{9. What to Record}

For every directed experiment, retain:

\begin{itemize}
\item workload file;
\item random seed, when applicable;
\item queue size;
\item exact executable used;
\item source revision or Git commit;
\item functional comparison result;
\item Multi-Push statistics;
\item characterization trace;
\item baseline statistics.
\end{itemize}

The characterization trace should be treated as the primary evidence for
the state-machine behavior being modeled.

\subsection{10. Interpretation Criteria}

The experiment is considered behaviorally consistent with the intended
Multi-Push mechanism when all of the following observations hold:

\begin{enumerate}
\item PUSH operations are accepted without immediately performing
HeapifyUp.

\item \texttt{push\_location} advances for every accepted PUSH.

\item \texttt{heap\_size} remains unchanged while pending maintenance
is outstanding.

\item A POP synchronizes all pending entries before removing the root.

\item After synchronization,
\(\texttt{push\_location}=\texttt{heap\_size}\).

\item After a POP, both state variables move consistently with the
reduced active heap.

\item The resulting POP sequence matches the conventional heap
implementation.


\end{enumerate}

This is a behavioral criterion, not a hardware-performance criterion.

\subsection{11. Current Experimental Conclusion}

The current software experiments support the following behavioral model:

PUSH acceptance→deferred HeapifyUp→POP synchronization→HeapifyDown

The directed traces demonstrate that
\texttt{push_location} and \texttt{heap_size} can be used to represent
the distinction between accepted entries and entries already incorporated
into the heap ordering.

The functional comparison against \texttt{heap_drain} also shows
identical POP output for the tested N=128 generated workload.

The appropriate next step is therefore not to modify the functional
reference model unnecessarily. The next engineering step is to use this
behavioral model to define the corresponding RTL state transitions and
then verify the RTL against the software reference.

\section{Key Terms}

\begin{description}[style=nextline,leftmargin=3.5cm]

\item[Multi-Push]
Optimization intended for workloads in which multiple PUSH operations
may arrive while heap maintenance is still in progress.

\item[PUSH acceptance]
The operation of storing an incoming entry at the current
\texttt{push\_location}. In the reference model, acceptance does not
immediately perform HeapifyUp.

\item[\texttt{push\_location}]
State variable identifying the next free tail position at which an
incoming PUSH can be stored.

\item[\texttt{heap\_size}]
State variable representing the number of entries currently
incorporated into the heap ordering.

\item[Pending PUSH]
An accepted entry that has been stored but has not yet completed its
HeapifyUp operation.

\item[HeapifyUp]
Heap operation that restores the heap property by moving a newly
incorporated entry toward the root when required.

\item[HeapifyDown]
Heap operation performed after root removal to restore the heap
property by moving the replacement entry toward the leaves when
required.

\item[Synchronization point]
A point at which the implementation must complete pending PUSH
maintenance before allowing a POP to remove the root.

\item[POP]
Priority-queue removal operation that returns the highest-priority
entry in the modeled min-heap.

\item[Characterization model]
A reference implementation whose purpose is to expose internal
algorithmic state transitions rather than to serve as the final
functional implementation.

\item[Functional reference]
The software model used as the behavioral source of truth for
comparison against other implementations.

\item[Baseline]
Conventional heap implementation used as an independent reference
for functional output and algorithmic operation counts.

\item[RTL]
Register-transfer level description of the hardware implementation.
The software experiment precedes RTL so that the expected state
transitions are defined before hardware design decisions are made.


\end{description}

\end{document}
